\chapter{Addiction Withdrawal}

\section*{Definition and Overview}
Withdrawal is the set of physical and mental symptoms that occur when a person who is dependent on a substance stops using it, or reduces their use, suddenly. It happens because the brain and body have adapted to the substance and ``rebound'' when it is removed.

Withdrawal can range from mild discomfort to a life-threatening medical emergency, depending on the substance, how long it was used, and the dose. Alcohol, benzodiazepines, opioids, nicotine, and stimulants each produce their own characteristic withdrawal syndrome, and the risks differ substantially between them.

\section*{Epidemiology}
Withdrawal affects people who have developed physical dependence, which is more common with regular, heavy, long-term use. Alcohol and benzodiazepines carry the greatest danger because their withdrawal can cause seizures and delirium.

Because many people try to stop these substances without medical help, withdrawal is a common reason for emergency department visits and for relapse: stopping is so unpleasant that many people use again simply to feel better. Medically supervised withdrawal markedly reduces the serious complications.

\section*{Aetiology and Pathophysiology}
Most addictive substances dampen or boost the nervous system, and the brain compensates over time to keep itself balanced. When the substance is removed, that compensation is unmasked. For example, alcohol slows brain activity, so the brain becomes more excitable in response; when drinking stops, this extra excitability appears as tremor, agitation, and, in severe cases, seizures. Opioids suppress pain pathways and the drive to breathe, so withdrawal produces the opposite effects, including pain, cramping, and agitation.

The severity of withdrawal depends on how strongly the substance has altered brain function, the size of the doses used, and how abruptly use stops. Tolerance and dependence develop together, and a rapid fall in tolerance during abstinence makes relapse especially dangerous.

\section*{Clinical Features}
Symptoms vary with the substance:

\begin{itemize}
    \item \textbf{Alcohol:} tremor, sweating, anxiety, nausea, poor sleep, and a raised heart rate; severe cases may progress to hallucinations, seizures, or delirium tremens
    \item \textbf{Benzodiazepines:} anxiety, insomnia, tremor, and panic; seizures can occur days to weeks after stopping
    \item \textbf{Opioids:} muscle aches, stomach cramps, diarrhoea, yawning, runny nose, sweating, and agitation
    \item \textbf{Nicotine:} cravings, irritability, poor concentration, low mood, and disturbed sleep
    \item \textbf{Stimulants:} fatigue, depression, increased appetite, disturbed sleep, and strong cravings
\end{itemize}

\section*{Differential Diagnosis}
Other conditions can produce similar symptoms, especially when the history of substance use is not known:

\begin{itemize}
    \item Hypoglycaemia (low blood sugar)
    \item Infection, such as sepsis or meningitis
    \item Stroke or head injury
    \item Thyroid disease or other hormonal problems
    \item Anxiety or panic disorder
    \item Medication side effects or interactions
\end{itemize}

\section*{When to Seek Medical Care}
Anyone withdrawing from alcohol or benzodiazepines, and anyone with serious symptoms or underlying illness, should be medically supervised. Even mild withdrawal deserves professional advice if there is any doubt, because symptoms can escalate rapidly.

\textbf{Call 000 (or local emergency services) for:}
\begin{itemize}
    \item Seizures or fits
    \item Confusion, disorientation, or hallucinations (possible delirium tremens)
    \item Severe agitation, fever, or a rapid heart rate
    \item Vomiting to the point of being unable to keep fluids down
    \item Suicidal thoughts or self-harm
\end{itemize}

\section*{Diagnosis and Investigations}
\begin{itemize}
    \item \textbf{History:} which substances are used, how much, for how long, and when the last dose was taken
    \item \textbf{Physical examination:} vital signs, tremor, agitation, and signs of dehydration
    \item \textbf{Withdrawal scales:} tools such as the CIWA scale help score the severity of alcohol withdrawal
    \item \textbf{Blood tests:} electrolytes, liver function, blood sugar, and alcohol level
    \item \textbf{ECG:} to check for heart rhythm problems
    \item \textbf{Mental state assessment:} to detect confusion, hallucinations, or risk of self-harm
\end{itemize}

\section*{Management}
Treatment depends on the substance and the severity of withdrawal:

\begin{itemize}
    \item \textbf{Supervised withdrawal (detox):} in hospital or a residential service for high-risk substances
    \item \textbf{Medicines:} such as benzodiazepines to prevent seizures in alcohol withdrawal, and symptom-relief medicines for opioid withdrawal
    \item \textbf{Substitution treatment:} longer-acting medicines, such as methadone or buprenorphine, for opioid dependence
    \item \textbf{Supportive care:} fluids, rest, nutrition, and a calm, quiet environment
    \item \textbf{Tapering:} gradual reduction rather than abrupt stopping, especially for benzodiazepines
    \item \textbf{Ongoing treatment:} counselling and relapse prevention, because withdrawal alone rarely cures addiction
\end{itemize}

\section*{Complications}
\begin{itemize}
    \item Seizures and delirium tremens (alcohol and benzodiazepines)
    \item Dehydration and electrolyte disturbances
    \item Heart rhythm problems
    \item Injury during seizures or confusion
    \item Severe depression and suicide risk
    \item Relapse, and overdose, because tolerance falls quickly during abstinence
\end{itemize}

\section*{Prognosis and Outcomes}
With medical supervision, most withdrawal symptoms settle within days to a few weeks. Alcohol withdrawal peaks in the first few days and usually improves within a week; opioid withdrawal is generally not life-threatening but can last longer. The dangerous complications are largely preventable when withdrawal is managed properly.

However, withdrawal is only the first step. Long-term recovery requires treatment of the underlying addiction, because successfully detoxifying without addressing the reasons for use leaves a very high risk of relapse.

\section*{Prevention and Self-Care}
\begin{itemize}
    \item Do not stop suddenly from alcohol or benzodiazepines without medical advice
    \item Seek professional help to plan withdrawal, especially for high-risk substances
    \item Have a support person present during withdrawal
    \item Stay hydrated and eat small, regular meals
    \item Avoid driving and operating machinery during withdrawal
    \item Arrange follow-up treatment before starting withdrawal
    \item Seek help quickly if symptoms become severe
\end{itemize}

\section*{Sources and Further Reading}
The following organisations provide reliable information on withdrawal and medically supervised detoxification.

\begin{itemize}
    \item NHS. \textit{Information on drug and alcohol withdrawal and detoxification.} https://www.nhs.uk/conditions/
    \item World Health Organization. \textit{Resources on substance use and withdrawal management.} https://www.who.int/
    \item NICE. \textit{Clinical guidance on alcohol dependence and detoxification.} https://www.nice.org.uk/
    \item MedlinePlus. \textit{Health information on drug and alcohol withdrawal.} https://medlineplus.gov/
    \item Mayo Clinic. \textit{Guide to withdrawal management and detoxification.} https://www.mayoclinic.org/
\end{itemize}
